%% kapa: exam=1 round=34 year=2023 period=2 form=A subject=appraisal_law subject_ko=감정평가관계법규 seq=20 number=20
%% source: Q-Net 감정평가사 기출문제 (cst00309) artlSeq=5211789 file=raw/1008/5211789/2023년도 제34회 감정평가사 1차시험 2교시.pdf
%% conversion: kordoc 4.24.51 -> tools/gen1.py (auto); review=proofread-sonnet; flags=none
\begin{question}{20}
국유재산법령상 행정재산과 일반재산에 관한 내용으로 옳지 않은 것은?
\begin{choices}
  \item 행정재산의 사용허가기간을 갱신받으려는 자는 허가기간이 끝나기 1개월 전에 중앙관서의 장에게 신청하여야 한다.
  \item 중앙관서의 장은 행정재산의 사용허가를 철회하려는 경우에는 청문을 하여야 한다.
  \item 일반재산의 대부계약은 수의계약의 방법으로 대부할 때에는 1회만 갱신할 수 있다.
  \item 행정재산의 사용허가가 취소된 경우에는 재산을 원래 상태대로 반환하여야 하지만, 중앙관서의 장이 미리 상태의 변경을 승인한 경우에는 변경된 상태로 반환할 수 있다.
  \item 일반재산을 매각한 이후 매수자가 매각대금을 체납한 경우 그 계약을 해제할 수 있다.
\end{choices}
\end{question}
\answer{3}
